% SOURCE_AUTHORITY: sga5_fr_workpass.tex, lines 8631--8647 (LF-normalized unit).
% SOURCE_SHA256: 791F4EFFC5E02832D5D77ED03518C8156D6F07E4C8238B03545DB93D883FBB28
\subsubsection*{1.4.4}
Si \(A\) es una \(\mathbb Z_\ell\)-álgebra conmutativa finita y \(F\) y \(G\) son dos \(A\)-haces constructibles, el sistema proyectivo
\[
F\otimes_A G=(F_n\otimes_{A_X} G_n)_{n\in\mathbb N},
\]
con los morfismos de transición evidentes, es un \(\mathbb Z_\ell\)-haz. Además, \(A\) actúa de manera evidente sobre cada una de sus componentes, lo que permite dotarlo de manera natural de una estructura de \(A\)-haz. Se define así un bifuntor exacto por la derecha, llamado producto tensorial,
\[
A\text{-}\fc(X)\times A\text{-}\fc(X)\longrightarrow A\text{-}\fc(X).
\]

Sean ahora \(F\) y \(G\) dos \(A\)-haces y supongamos que \(F\) es localmente libre, es decir, que, para todo entero \(n\), \(F_n\) es un \((A/\ell^{\,n+1}A)_X\)-módulo localmente libre. Entonces, al igual que en (1.3.2), se ve que el sistema proyectivo
\[
\cHom_A(F,G)=\bigl(\cHom_{A_n}(F_n,G_n)\bigr)_{n\in\mathbb N},
\]
con los morfismos de transición evidentes, es un \(\mathbb Z_\ell\)-haz. Además, como \(A\) actúa naturalmente sobre sus componentes, está dotado canónicamente de una estructura de \(A\)-haz.

Al igual que en (1.3.3), se define la noción de \(A\)-haz invertible y, si \(L\) es un tal \(A\)-haz, se da sentido a la expresión \(L^{\otimes r}\) para todo entero natural \(r\). El apartado (1.3.3) se adapta sin dificultad a esta situación.
